\section{Program representation}

\subsection{Flow diagram representation}
The rules for the representation of a flow diagram are the following,
right after, a familiar diagram is shown.
\begin{itemize}
\item Circles indicate insertion and termination points,
  respectively named ``start'' and ``end''.
\item Slanted parallelograms indicate insertion variables
  also known as \textbf{input} and termination variables also
  know as \textbf{output}.
\item Rhombi are \textbf{conditionals} their \underline{exit flow is not fixed},
  its a function of their inner state.
  For our purposes the possible exit flows are \textit{two}, one if the condition
  \underline{is} satisfied and another one if it \underline{is not}.
\item Rectangles are \textbf{procedures}, their \underline{exit flow is fixed}.
\end{itemize}

\begin{figure}[h]
  \centering
  \includegraphics[width=\textwidth]{flow-diagram-example.png}
  \caption{
    This is the flow diagram of the ``factorial'' program
    previously shown.
  }
  \label{fig:flow-diagram-example}
\end{figure}

\pagebreak

\subsection{Procedural representation}
The ordering given by arrows is now replaced 
There are no more arrows, each line is executed in order
from top to bottom up until an \textit{indirection} command
such as ``GOTO'' which diverts the next execution to a
special line marked with a ``FLAG'' statement. This is a
typical scheme for how computers execute code.
\begin{minted}[bgcolor=mygreen!10]{text}
  START n
  SETTO x 1
  FLAG LOOP
  IFSTART GREATERTHAN n 0
  SETTO x PRODUCT x n
  SETTO n SUBTRACTION n 1
  GOTO LOOP
  IFSTOP
  ELSESTART
  STOP x
  ELSESTOP
\end{minted}

\pagebreak

\subsection{Flow patterns}
The flow of a program is a very important aspect to understand its behaviour,
all useful programs have a source-sink aspect to them, all arrows begin from
a source and they must all end in a sink at some point.

\noindent
\begin{minipage}{0.5\textwidth}
  \vspace{1cm}
  \includegraphics[width=\textwidth]{flow-if-else.png}
  \captionof{figure}{This is an \textbf{if-else} pattern, it starts from $C$.}
  \label{fig:flow-if-else}
\end{minipage}
\hspace{0.05\textwidth}
\begin{minipage}{0.4\textwidth}
  The \textit{if-else} pattern \underline{conserves} the tree structure of the program.
  If $1$ and $2$ are halting programs than the whole tree will halt.
  \begin{minted}[bgcolor=mygreen!10]{c}
    if (C) {
      // 1
    }
    else {
      // 2
    }
  \end{minted}
\end{minipage}

\noindent
\begin{minipage}{0.4\textwidth}
  \vspace{1cm}
  The \textit{if-elseif} pattern is a pattern extension of the \textit{if-else} pattern.
  It \underline{conserves} the tree structure of the program.
  If $1 \cdots n$ are halting programs than the whole tree will halt.
  \begin{minted}[bgcolor=mygreen!10]{c}
    if (1) {
      // 1
    }
    else if (...) {
      // ...
    }
    else {
      // N
    }
  \end{minted}
\end{minipage}
\hspace{0.05\textwidth}
\begin{minipage}{0.5\textwidth}
  \includegraphics[width=\textwidth]{flow-if-elseif.png}
  \captionof{figure}{This is an \textbf{if-elseif} pattern, it starts from $C$.}
  \label{fig:flow-if-elseif}
\end{minipage}

\noindent
\begin{minipage}{0.5\textwidth}
  \vspace{1cm}
  \includegraphics[width=\textwidth]{flow-if.png}
  \captionof{figure}{This is a free \textbf{if} pattern, it starts from $C$.}
  \label{fig:flow-if}
\end{minipage}
\hspace{0.05\textwidth}
\begin{minipage}{0.4\textwidth}
  The \textit{free if} pattern is a \underline{reduction} of an \textit{if-else} pattern.
  If $1 \otimes 2$ and $2$ halts than the whole tree will halt
  \begin{minted}[bgcolor=mygreen!10]{c}
    if (C) {
      // 1
    }
    // 2
  \end{minted}
\end{minipage}

\noindent
\begin{minipage}{0.4\textwidth}
  \vspace{1cm}
  The \textit{while} pattern is guaranteed to halt
  if the condition $C$ has an \textit{upper-bound}
  check as a logical conjunction and is a function
  of a \textit{strictly increasing} value (the same
  is true for a lower-bound with a stricly decreasing
  value).
  \begin{minted}[bgcolor=mygreen!10]{c}
    while (C) {
      // B
    }
  \end{minted}
\end{minipage}
\hspace{0.05\textwidth}
\begin{minipage}{0.5\textwidth}
  \includegraphics[width=\textwidth]{flow-while.png}
  \captionof{figure}{This is \textbf{while} pattern, it starts from $C$.}
  \label{fig:flow-while}
\end{minipage}

\noindent
\begin{minipage}{0.5\textwidth}
  \vspace{1cm}
  \includegraphics[width=\textwidth]{flow-do-while.png}
  \captionof{figure}{This is a \textbf{do-while} pattern, it starts from $C$.}
  \label{fig:flow-do-while}
\end{minipage}
\hspace{0.05\textwidth}
\begin{minipage}{0.4\textwidth}
  The \textit{do-while} pattern exhibits the same
  halting behaviour as the \textit{while} pattern;
  the difference is that the \textit{body} of the
  loop $B$ will always resolve \underline{at least}
  once.
  \begin{minted}[bgcolor=mygreen!10]{c}
    do {
      // B
    } while (C);
  \end{minted}
\end{minipage}

\noindent
\begin{minipage}{0.4\textwidth}
  \vspace{1cm}
  The \textit{for} pattern exhibits the same
  halting behaviour as the \textit{while} pattern;
  in addition if $C \otimes B = C \oplus B$ then
  the strictly increasing/decreasing value is
  controlled by $IT$.
  \begin{minted}[bgcolor=mygreen!10]{c}
    for (IN; C; IT) {
      // B
    }
  \end{minted}
\end{minipage}
\hspace{0.05\textwidth}
\begin{minipage}{0.5\textwidth}
  \includegraphics[width=\textwidth]{flow-for.png}
  \captionof{figure}{This is \textbf{for} pattern, it starts from $IN$.}
  \label{fig:flow-for}
\end{minipage}

\paragraph{Reductions}
